\chapter{Betinget sandsynlighed og Bayes' sætning}
\label{chap:ch5}

I dette kapitel introducerer vi mere avancerede begreber inden for sandsynlighedsregning. Vi begynder med \textbf{betinget sandsynlighed}, som vurderer sandsynligheden for, at en hændelse indtræffer, \textit{givet} at en anden hændelse allerede har fundet sted. I forlængelse heraf introducerer vi \textbf{multiplikationsreglen}, et centralt princip til at bestemme sandsynligheden for, at flere hændelser indtræffer efter hinanden. Dernæst ser vi på \textbf{loven om total sandsynlighed}, som gør det muligt at opdele og beregne sandsynligheder på tværs af forskellige scenarier eller partitioner af udfaldsrummet. For bedre at kunne skelne mellem, hvordan hændelser påvirker hinanden, undersøger vi \textbf{afhængige og uafhængige hændelser}. Her afklarer vi, hvordan forekomsten af én hændelse påvirker eller ikke påvirker en anden.

Til sidst ser vi nærmere på \textbf{Bayes' sætning}, et stærkt værktøj til at opdatere sandsynligheder i lyset af ny evidens.

\section{Betinget sandsynlighed}
Se på følgende scenarie: I en bestemt population har 5\% af individerne en bestemt sygdom. Sandsynligheden for, at en tilfældigt udvalgt person har sygdommen, er derfor 5\%:

\[
P(D) = 0.05,
\]

Her betegner \( D \) hændelsen, at personen har sygdommen. \( P(D) \) kaldes \textit{a priori-sandsynligheden}: vores vurdering før yderligere information.

Forestil dig nu, at vi udvælger en tilfældig person og får at vide, at personen er testet positiv for sygdommen. Hvordan skal vi med denne ekstra information opdatere sandsynligheden for, at personen faktisk har sygdommen? Hvad er med andre ord sandsynligheden for, at en person har sygdommen, givet at personen er testet positiv? Lad \( T \) betegne hændelsen, at en person testes positiv. Denne betingede sandsynlighed skrives som:

\[
P(D \mid T),
\]

hvilket er sandsynligheden for \( D \), givet at \( T \) er indtruffet. Den opdaterede sandsynlighed \( P(D \mid T) \) kaldes \textit{a posteriori-sandsynligheden}, fordi den udtrykker vores reviderede opfattelse af hændelsens sandsynlighed, efter at vi har taget den nye evidens i betragtning. Intuitivt er det rimeligt at forvente, at \( P(D \mid T) \) er større end a priori-sandsynligheden \( P(D) \). Men hvad er den præcise værdi af \( P(D \mid T) \)? Før vi introducerer en generel formel, ser vi på et simpelt eksempel.

\begin{example}
Jeg kaster en fair terning. Lad \( A \) være hændelsen, at udfaldet er et primtal, dvs. \( A = \{2, 3, 5\} \). Lad desuden \( B \) være hændelsen, at udfaldet er større end eller lig med 3, dvs. \( B = \{3, 4, 5, 6\} \).

\begin{enumerate}[label=(\alph*)]
    \item Hvad er sandsynligheden for \( A \), \( P(A) \)?
    \item Hvad er sandsynligheden for \( A \) givet \( B \), \( P(A \mid B) \)?
\end{enumerate}

\end{example}

\begin{solution}
    \begin{enumerate}[label=(\alph*)]
        \item Sandsynligheden for \( A \) er antallet af udfald i \( A \) divideret med det samlede antal udfald. Da der er 3 primtal og 6 mulige udfald, har vi:
    
        \begin{align*}
            P(A) &= \frac{3}{6} = \frac{1}{2}
        \end{align*}
    
        \item Sandsynligheden for \( A \) givet \( B \) er antallet af udfald i fællesmængden af \( A \) og \( B \) divideret med antallet af udfald i \( B \). Da udfaldene i fællesmængden af \( A \) og \( B \) er \(\{3, 5\}\), og udfaldene i \( B \) er \(\{3, 4, 5, 6\}\), har vi: \vspace{0.4cm}
    
        \begin{align*}
        P(A \mid B) &=\frac{|A \cap B|}{|B|} = \frac{|\{3,5\}|}{|\{3,4, 5, 6\}|} = \frac{1}{2}
        \end{align*}

    \end{enumerate}
\end{solution}

Nu hvor vi har forstået det grundlæggende eksempel, kan vi generalisere fremgangsmåden og udlede en mere generel formel for betinget sandsynlighed. Med udgangspunkt i det konkrete tilfælde kan vi omskrive udtrykket ved at dividere både tæller og nævner med det samlede antal mulige udfald, \( |S| \), som vist nedenfor:

\[
P(A \mid B) = \frac{|A \cap B|}{|B|} = \frac{\frac{|A \cap B|}{|S|}}{\frac{|B|}{|S|}} = \frac{P(A \cap B)}{P(B)}
\]

\textbf{Forklaring:}

\begin{itemize}
    \item \textbf{Tæller:} \( |A \cap B| \) er antallet af udfald, hvor både hændelse \( A \) og hændelse \( B \) indtræffer. Ved at dividere med \( |S| \) omdannes dette antal til en sandsynlighed, \( P(A \cap B) \).
    \item \textbf{Nævner:} \( |B| \) er antallet af udfald, hvor hændelse \( B \) indtræffer. Ved tilsvarende at dividere med \( |S| \) får vi sandsynligheden for hændelse \( B \), som betegnes \( P(B) \).
\end{itemize}

Den betingede sandsynlighed \( P(A \mid B) \) kan altså udtrykkes som forholdet mellem den fælles sandsynlighed for \( A \) og \( B \) og sandsynligheden for \( B \).

Selv om udledningen ovenfor forudsætter et endeligt udfaldsrum med lige sandsynlige udfald, er den resulterende formel bemærkelsesværdigt generel. Den gælder, uanset om udfaldsrummet er endeligt eller uendeligt, og uanset om udfaldene er lige sandsynlige eller ej. Denne generalitet gør formlen til en hjørnesten i sandsynlighedsteorien.

\begin{definition}[Betinget sandsynlighed]
For to vilkårlige hændelser \( A \) og \( B \) med \( P(B) > 0 \) er den betingede sandsynlighed for \( A \) givet \( B \) defineret som:
\[
P(A \mid B) = \frac{P(A \cap B)}{P(B)}, \qquad P(B) > 0
\]
\end{definition}

Her er intuitionen bag formlen. Når vi ved, at hændelse \( B \) er indtruffet, udelukker vi reelt alle udfald, der ikke er en del af \( B \). Derved bliver vores oprindelige udfaldsrum reduceret til delmængden \( B \).

I dette nye, begrænsede udfaldsrum \( B \) kan hændelse \( A \) kun indtræffe, hvis udfaldet ligger i fællesmængden af \( A \) og \( B \), som betegnes \( A \cap B \). For at bestemme den betingede sandsynlighed \( P(A \mid B) \) beregner vi forholdet mellem sandsynligheden for, at både \( A \) og \( B \) indtræffer, og sandsynligheden for, at \( B \) indtræffer alene. Matematisk udtrykkes det som:

\[
P(A \mid B) = \frac{P(A \cap B)}{P(B)}
\]

Denne formulering sikrer, at sandsynlighederne i det nye udfaldsrum \( B \) er normaliserede, dvs. at den samlede sandsynlighed summer til 1. Se for eksempel på den betingede sandsynlighed for \( B \) givet \( B \):

\[
P(B \mid B) = \frac{P(B \cap B)}{P(B)} = \frac{P(B)}{P(B)} = 1
\]

Dette resultat bekræfter intuitivt, at hvis \( B \) er indtruffet, er det sikkert, at \( B \) indtræffer, dvs. sandsynligheden er 1.

Bemærk, at den elementære formel for \( P(A \mid B) \) ikke er defineret, når \( P(B) = 0 \), fordi den ville kræve division med nul. En hændelse med sandsynlighed nul behøver ikke at være tom i en kontinuert sandsynlighedsmodel, men betingning på sådanne hændelser kræver mere avancerede definitioner, som ligger uden for dette kapitels rammer.

Vi kan også formulere sandsynlighedsregningens aksiomer ved hjælp af betinget sandsynlighed. Aksiomerne er som følger:

\begin{axiom}[Aksiomer for betinget sandsynlighed]
    \begin{itemize}
        \item \textbf{Aksiom 1:} For enhver hændelse \(A\) gælder \( 0 \leq P(A\mid B) \leq 1 \).
        \item \textbf{Aksiom 2:} Den betingede sandsynlighed for $B$ givet $B$ er 1, dvs. $P(B \mid B)=1$.
        \item \textbf{Aksiom 3:} Hvis \(A_1, A_2, A_3, \cdots\) er disjunkte hændelser, så gælder
         
        \hspace*{2.8em}$P\left(A_1 \cup A_2 \cup A_3 \cdots \mid B\right)=P\left(A_1 \mid B\right)+P\left(A_2 \mid B\right)+P\left(A_3 \mid B\right)+\cdots$
    \end{itemize}
\end{axiom}

Ud fra disse aksiomer kan vi også udlede andre regler for betinget sandsynlighed:

\begin{theorem}{Regler for betinget sandsynlighed}
    \begin{itemize}
        \item \textbf{Komplementreglen:} Sandsynligheden for komplementet til hændelse $A$ givet $C$ er
    
        \[
        P\left(A^c \mid C\right)=1-P(A \mid C)
        \]

        \item \textbf{Reglen for den tomme mængde:} Sandsynligheden for den tomme mængde givet en hændelse $C$ er 0, dvs.
        
        \[P(\emptyset \mid C) = 0\]
        
        \item \textbf{Additionsreglen:} For to vilkårlige hændelser $A$ og $B$ er sandsynligheden for foreningsmængden af hændelserne $A$ og $B$ givet en anden hændelse $C$
        \[
        P(A \cup B \mid C)=P(A \mid C)+P(B \mid C)-P(A \cap B \mid C)
        \]

        \item \textbf{Differensreglen:} Sandsynligheden for differensen mellem hændelserne $A$ og $B$ givet en anden hændelse $C$ er
        \[
            P(A-B \mid C)=P(A \mid C)-P(A \cap B \mid C)
        \]
        \item \textbf{Delmængdereglen:} Hvis $A \subset B$, så gælder
        \[
        P(A \mid C) \leq P(B \mid C)
        \]
    \end{itemize}
\end{theorem}

Det følgende eksempel viser, hvordan betinget sandsynlighed anvendes, og introducerer begrebet \textbf{kontingenstabeller} samt \textbf{falsk positive} og \textbf{falsk negative}.

\begin{example}
    En forsker vil vurdere, hvor effektiv en diagnostisk test er til at påvise nyresygdom hos patienter med forhøjet blodtryk. Hun udfører derfor testen på en stikprøve af 137 patienter, som er inddelt således:

\begin{itemize}
    \item \textbf{67 patienter} med en bekræftet diagnose af nyresygdom.
    \item \textbf{70 patienter}, som vides at være raske (dvs. ikke har nyresygdom).
\end{itemize}

Den diagnostiske test giver et af to mulige resultater for hver patient:

\begin{itemize}
    \item \textbf{Positiv}: Angiver, at patienten har nyresygdom.
    \item \textbf{Negativ}: Angiver, at patienten ikke har nyresygdom.
\end{itemize}

Resultaterne er opsummeret i følgende kontingenstabel:

    \begin{tabular}{lrrr} 
        & \multicolumn{2}{c}{ Testresultat } & \\
        Faktisk tilstand & Positiv & Negativ & I alt \\
        Nyresygdom & 44 & 23 & 67 \\
        Rask & 10 & 60 & 70 \\
        I alt & 54 & 83 & 137
        \end{tabular}
        
    I dette eksperiment er:


        \begin{itemize}
            \item \textbf{Sandt positive}: Patienter, der har nyresygdom og er testet positive.
            \item \textbf{Falsk negative}: Patienter, der har nyresygdom, men er testet negative.
            \item \textbf{Falsk positive}: Raske patienter, der er testet positive.
            \item \textbf{Sandt negative}: Raske patienter, der er testet negative.
        \end{itemize}
    
        Bestem følgende sandsynligheder:
        \begin{enumerate}[label=(\alph*)]
            \item Sandsynligheden for at have nyresygdom, $P(D)$.
            \item Sandsynligheden for en positiv test,  \( P(T^+) \).
            \item Sandsynligheden for en negativ test,  \( P(T^-) \).
            \item Hvis en person har nyresygdom, hvad er så sandsynligheden for, at personen testes positiv for sygdommen?
            \item Bestem sandsynligheden for, at en patient har nyresygdom givet et positivt testresultat, \( P(D \mid T^+) \).
            \item Bestem sandsynligheden for, at en patient ikke har nyresygdom givet et negativt testresultat, \( P(\text{Rask} \mid T^-) \).
            \item Vurdér den diagnostiske tests samlede nøjagtighed.
        \end{enumerate}
        \label{ex:renal-disease}
\end{example}

\begin{solution}
    Ud fra kontingenstabellen kan vi beregne følgende sandsynligheder:

    \begin{enumerate}[label=(\alph*)]
        \item \textbf{Sandsynligheden for nyresygdom, \( P(D) \)}:
        \begin{align*}
        P(D) &= \frac{\text{Antal patienter med nyresygdom}}{\text{Samlet antal patienter}} \\
        &= \frac{67}{137} \approx 0.4891 \text{ eller } 48.91\%
        \end{align*}
        
        \item \textbf{Sandsynligheden for en positiv test, \( P(T^+) \)}:
        \begin{align*}
        P(T^+) &= \frac{\text{Antal positive tests}}{\text{Samlet antal patienter}} \\
        &= \frac{54}{137} \approx 0.3942 \text{ eller } 39.42\%
        \end{align*}
        
        \item \textbf{Sandsynligheden for en negativ test, \( P(T^-) \)}:
        \begin{align*}
        P(T^-) &= \frac{\text{Antal negative tests}}{\text{Samlet antal patienter}} \\
        &= \frac{83}{137} \approx 0.6058 \text{ eller } 60.58\%
        \end{align*}


        \item \textbf{Sandsynligheden for en positiv test givet nyresygdom, \( P(T^+ \mid D) \)}:
        \begin{align*}
            P(T^+ \mid D)=\frac{44}{67} \approx 0.6567
        \end{align*}

        \item \textbf{Sandsynligheden for nyresygdom givet en positiv test, \( P(D \mid T^+) \)}:
        \begin{align*}
        P(D \mid T^+) &= \frac{P(D \cap T^+)}{P(T^+)} \\
        &= \frac{44}{54}=\frac{22}{27} \approx 0.8148
        \end{align*}

        \item \textbf{Sandsynligheden for at være rask givet en negativ test, \( P(\text{Rask} \mid T^-) \)}:
        \begin{align*}
        P(\text{Rask} \mid T^-) &= \frac{P(\text{Rask} \cap T^-)}{P(T^-)} \\
        &= \frac{60}{83} \approx 0.7229
        \end{align*}

        \item \textbf{Den diagnostiske tests samlede nøjagtighed}:
        Den diagnostiske tests samlede nøjagtighed er andelen af korrekte diagnoser, dvs. summen af sandt positive og sandt negative divideret med det samlede antal patienter:
        \begin{align*}
            \text{Samlet nøjagtighed} =\frac{44+60}{137}=\frac{104}{137} \approx 0.7591
        \end{align*}
    \end{enumerate}


Beregningerne beskriver flere forskellige mål for testens ydeevne. Sensitiviteten, som er sandsynligheden for et positivt resultat blandt patienter med nyresygdom, er \(44/67\approx65.67\%\). Specificiteten, som er sandsynligheden for et negativt resultat blandt raske patienter, er \(60/70\approx85.71\%\). Derimod er \(P(D\mid T^+)=81.48\%\) den positive prædiktive værdi, og \(P(\text{Rask}\mid T^-)=72.29\%\) er den negative prædiktive værdi. Den samlede nøjagtighed er 75.91\%.
\end{solution}

Her er et klassisk sandsynlighedsproblem, som kaldes tobørnsproblemet. Scenariet optræder i forskellige udgaver i litteraturen, og hver udgave har små drejninger, som fører til forskellige resultater. Før du ser på udregningerne, så brug et øjeblik på at komme med dine egne forudsigelser -- du bliver måske overrasket over resultaterne!

\begin{example}
    Se på en familie med to børn. Vi er interesserede i de mulige kombinationer af børnenes biologiske køn. Udfaldsrummet for denne situation er:

\[
S = \{(G, G), (G, B), (B, G), (B, B)\}
\]

hvor \( G \) betegner en biologisk pige (af engelsk »girl«), og \( B \) betegner en biologisk dreng (af engelsk »boy«). For enkelhedens skyld antager vi, at alle fire udfald er lige sandsynlige.

\begin{enumerate}[label=(\alph*)]
    \item Hvad er sandsynligheden for, at begge børn er piger, givet at det første barn er en pige?
    \item Vi spørger faren: »Har du mindst én datter?« Han svarer »Ja!« Hvad er sandsynligheden for, at begge børn er piger, givet denne ekstra information? Hvad er med andre ord sandsynligheden for, at begge børn er piger, givet at vi ved, at mindst ét af dem er en pige?
\end{enumerate}
\end{example}

\begin{solution}
    Lad $A$ være hændelsen, at begge børn er piger, dvs. $A=\{(G, G)\}$. Lad $B$ være hændelsen, at det første barn er en pige, dvs. $B=\{(G, G),(G, B)\}$. Lad til sidst $C$ være hændelsen, at mindst ét af børnene er en pige, dvs. $C=\{(G, G),(G, B),(B, G)\}$. Da udfaldene er lige sandsynlige, kan vi skrive
    \begin{align*}
    P(A) &= \frac{1}{4} \\
    P(B) &= \frac{2}{4} = \frac{1}{2} \\
    P(C) &= \frac{3}{4}
    \end{align*}


    \begin{enumerate}[label=(\alph*)]
         \item Hvad er sandsynligheden for, at begge børn er piger, givet at det første barn er en pige? Det er $P(A \mid B)$, så vi kan skrive

                    \begin{align*}
                    P(A \mid B) & =\frac{P(A \cap B)}{P(B)} \\
                    & =\frac{P(A)}{P(B)} \quad(\text { da } A \subset B) \\
                    & =\frac{\frac{1}{4}}{\frac{1}{2}}=\frac{1}{2}
                    \end{align*}
                
                \item Hvad er sandsynligheden for, at begge børn er piger, givet at vi ved, at mindst ét af dem er en pige? Det er $P(A \mid C)$, så vi kan skrive

                    \begin{align*}
                    P(A \mid C) & =\frac{P(A \cap C)}{P(C)} \\
                    & =\frac{P(A)}{P(C)} \quad(\text { da } A \subset C) \\
                    & =\frac{\frac{1}{4}}{\frac{3}{4}}=\frac{1}{3}
                    \end{align*}
        
    \end{enumerate}
\end{solution}

\begin{remark}
    Mange vil intuitivt gætte på, at både \( P(A \mid B) \) og \( P(A \mid C) \) er 50\%. Men selv om \( P(A \mid B) = 50\%\), er \( P(A \mid C) \) kun 33\%. Det viser, hvor kontraintuitiv sandsynlighed kan være. Nøglen er at indse, at hændelse \( B \) er en delmængde af hændelse \( C \). Konkret udelukker \( B \) udfaldet \( (B, G) \), som er med i \( C \). Derfor har \( C \) flere udfald, der ikke ligger i \( A \), end \( B \) har, hvilket giver en mindre \( P(A \mid C) \) sammenlignet med \( P(A \mid B) \).

    Resultatet i del (b) forudsætter, at den givne information er præcis hændelsen \(C\), »mindst ét barn er en pige«. Hvis faren blev udvalgt eller spurgt efter en anden tilfældig procedure, kan den betingede sandsynlighed ændre sig. Når man fortolker dette problem, skal man derfor angive den procedure, der frembringer informationen.
\end{remark}

\section{Multiplikationsreglen og loven om total sandsynlighed}
Multiplikationsreglen er et grundlæggende princip i sandsynlighedsteorien, som gør det muligt at beregne sandsynligheden for, at flere hændelser indtræffer efter hinanden. Reglen er særlig nyttig, når hændelserne er afhængige, dvs. når forekomsten af én hændelse påvirker sandsynligheden for den efterfølgende hændelse.

\begin{theorem}[Multiplikationsreglen]
    Hvis \(P(A)>0\) og \(P(B)>0\), er sandsynligheden for, at begge hændelser indtræffer, givet ved:
    \[
    P(A \cap B) = P(A) P(B \mid A) = P(B) P(A \mid B)
    \]
\end{theorem}

\begin{example}
    Antag, at et softwareudviklingsteam tester en ny funktion. Testen foregår i to trin: enhedstest (A) og integrationstest (B). Sandsynligheden for, at en fejl opdages under enhedstesten, er $P(A)=0.3$. Hvis en fejl opdages under enhedstesten, er sandsynligheden for, at den også opdages under integrationstesten, $P(B \mid A) = 0.7$. Hvad er sandsynligheden for, at en fejl opdages i begge testtrin?
\end{example}

\begin{solution}
    Med multiplikationsreglen får vi:

\[
P(A \cap B)=P(A) \cdot P(B \mid A)=0.3 \times 0.7=0.21
\]


Sandsynligheden for, at en fejl opdages både under enhedstesten og integrationstesten, er altså 0.21.
\end{solution}

\begin{example}
En virksomhed har indført et sikkerhedssystem i to lag for at beskytte mod netværksindbrud: \textbf{firewall (A)} og \textbf{indtrængningsdetektionssystem (IDS) (B)}. Sandsynligheden for, at et indbrudsforsøg opdages af firewallen, er \( P(A) = 0.4 \). Hvis indbruddet slipper igennem firewallen, er sandsynligheden for, at det opdages af IDS'en, \( P(B \mid A^c) = 0.6 \). Givet at firewallen opdager indbruddet, er den betingede sandsynlighed for, at IDS'en også opdager det, \( P(B \mid A) = 0.8 \).

\begin{enumerate}[label=(\alph*)]
    \item Hvad er sandsynligheden for, at et indbrud opdages af mindst ét af sikkerhedslagene?
    \item Hvad er sandsynligheden for, at et indbrud opdages af begge sikkerhedslag?
\end{enumerate}

\end{example}

\begin{solution}
\begin{enumerate}[label=(\alph*)]
    \item \textbf{Sandsynligheden for at blive opdaget af mindst ét lag:}

    For at finde sandsynligheden for, at indbruddet opdages af mindst ét lag, beregner vi sandsynligheden for, at et indbrud slipper uopdaget gennem begge lag, og trækker den fra 1.

    \begin{align*}
        P(\text{Uopdaget}) &= P(A^c) \cdot P(B^c \mid A^c) 
        = (1 - 0.4) \cdot (1 - 0.6) 
        = 0.6 \cdot 0.4
        = 0.24
    \end{align*}

    Sandsynligheden for at opdage indbruddet med mindst ét lag er derfor:

    \begin{align*}
    P(\text{Opdaget}) &= 1 - P(\text{Uopdaget}) = 1 - 0.24 = 0.76
    \end{align*}

    \item \textbf{Sandsynligheden for at blive opdaget af begge lag:}

    Med multiplikationsreglen får vi:

    \begin{align*}
    P(A \cap B) &= P(A) \cdot P(B \mid A) = 0.4 \cdot 0.8 = 0.32
    \end{align*}

    Sandsynligheden for at opdage indbruddet med mindst ét lag er 0.76, mens sandsynligheden for at opdage indbruddet med både firewallen og IDS'en er 0.32. Denne forskel understreger, hvor vigtigt det er at tage højde for afhængigheder mellem hændelser, når man beregner sandsynligheder.
\end{enumerate}
\end{solution}

\begin{remark}
    Man kunne fejlagtigt antage, at de to lag er uafhængige, og beregne $P(A \cap B)=P(A) \cdot P(B)=0.4 \cdot 0.68=0.272$, hvor $P(B)=0.4 \cdot 0.8+0.6 \cdot 0.6=0.68$. Da sandsynligheden for, at IDS'en opdager indbruddet, afhænger af firewallens resultat, bruger vi i stedet $P(B \mid A)$ og får det korrekte resultat 0.32.
\end{remark}

I nogle situationer afhænger sandsynligheden for en hændelse af forskellige betingelser. Hvis vi kender de betingede sandsynligheder under disse forskellige scenarier, kan vi bestemme hændelsens samlede sandsynlighed. Se for eksempel på halvlederproduktion, hvor vi definerer:
\begin{itemize}
    \item \( A \) som hændelsen, at en chip er stærkt forurenet
    \item \( B \) som hændelsen, at et produkt, der bruger chippen, svigter
\end{itemize}
Sandsynligheden for svigt for en chip, der ikke er forurenet, er \( P(B \mid A^c) = 0.005 \). Hvis chippen derimod er udsat for et højt forureningsniveau, er sandsynligheden for svigt \( P(B \mid A) = 0.10 \). Antag, at \( 20\% \) af chippene i en given produktionsserie er stærkt forurenede (\( P(A) = 0.20 \)). Hvad er sandsynligheden for, at et produkt, der bruger en af disse chips, svigter?

Sandsynligheden for svigt afhænger af, om chippen har været udsat for høj forurening eller ej. Enhver hændelse \( A \) kan opdeles i to gensidigt udelukkende dele: én, der skærer \( B \), og én, der skærer komplementet \( B^c \). Matematisk kan vi udtrykke det som:

\[
A = (B \cap A) \cup (B^c \cap A)
\]

Denne opdeling er vist i Venn-diagrammet i Figur \ref{fig:decoA}. Da \( B \) og \( B^c \) er gensidigt udelukkende, er deres fællesmængder med \( A \) også gensidigt udelukkende. Ved at bruge reglen for sandsynligheden for foreningsmængden af gensidigt udelukkende hændelser og multiplikationsreglen kan vi derfor udlede den totale sandsynlighed således:

\[
P(A) = P(A \cap B) + P(A \cap B^c)
\]

% Definition of ellipses
\def\firstellipse{(3,1) ellipse (4.5cm and 2.2cm)} % First ellipse with horizontal and vertical radii

\colorlet{ellipse edge}{main}
\colorlet{ellipse area}{main!20}

\tikzset{filled/.style={fill=ellipse area, draw=ellipse edge, thick},
    outline/.style={draw=ellipse edge, thick}}

\colorlet{ellipse edge}{main}
\colorlet{ellipse area}{main!20}

\begin{figure}[htbp]
    \centering
    \begin{tikzpicture}
        % Draw the ellipses

        % Draw the rectangle to represent the universal set
        \draw[outline] (-2,-2) rectangle (8,4);
        

         % Draw the ellipse and fill
        \begin{scope}
            \fill[filled] \firstellipse;
            \end{scope}
            \draw[outline] \firstellipse node[anchor=north west, xshift=2.5cm, yshift=-1.8cm] {$A$} node[anchor=north west, xshift=1.5cm, yshift=0.2cm] {$A \cap B^c$} node[anchor=north west, xshift=-4.2cm, yshift=2.8cm] {$B$} node[anchor=north west, xshift=-3cm, yshift=0.2cm] {$A \cap B$} node[anchor=north west, xshift=3.7cm, yshift=2.8cm] {$B^c$};

         % Draw a squiggly line down the middle of the rectangle
         \draw[decorate, decoration={snake, amplitude=0.05cm, segment length=1cm, post length=0.5cm}] (3,4) -- (3,-2);
        \end{tikzpicture}
        \caption{$P(A)=P(A \cap B)+P(A \cap B^c)$}
        \label{fig:decoA}
\end{figure}

Vi opsummerer dette i følgende sætning:

\begin{theorem}[Loven om total sandsynlighed]
For vilkårlige hændelser \( A \) og \( B \), hvor \( B \) og \( B^c \) udgør en partition af udfaldsrummet \( S \), er den totale sandsynlighed for hændelse \( A \) givet ved:
    \begin{align*}
    P(A) &= P(A \cap B) + P(A \cap B^c) \\
    &= P(B) P(A \mid B) + P(B^c) P(A \mid B^c)
    \end{align*}

    For enhver hændelse \( A \) og enhver partition \( B_1, B_2, \ldots, B_n \) af udfaldsrummet, hvor \( B_i \cap B_j = \emptyset \) for alle \( i \neq j \) og \( \bigcup_{i=1}^{n} B_i = S \), er den totale sandsynlighed for hændelse \( A \) givet ved:
    \begin{align*}
    P(A) &= \sum_{i=1}^{n} P(A \mid B_i) P(B_i)
    \end{align*}
\end{theorem}

Lad os se på nogle eksempler.

\begin{example}
    En virksomhed producerer to typer produkter: \textbf{produkt A} og \textbf{produkt B}. Sandsynligheden for, at et produkt er defekt, er \( P(D \mid A) = 0.05 \) for produkt A og \( P(D \mid B) = 0.10 \) for produkt B. Virksomheden fremstiller 60\% af sine produkter som produkt A og 40\% som produkt B. Hvad er sandsynligheden for, at et tilfældigt udvalgt produkt er defekt?
\end{example}

\begin{solution}
    Med loven om total sandsynlighed får vi:

    \[
    P(D) = P(D \mid A) P(A) + P(D \mid B) P(B)
    \]

    Vi indsætter de givne værdier:

    \[
    P(D) = 0.05 \times 0.60 + 0.10 \times 0.40 = 0.03 + 0.04 = 0.07
    \]

    Sandsynligheden for, at et tilfældigt udvalgt produkt er defekt, er altså 0.07.
\end{solution}

En grafisk fremstilling af, hvordan en hændelse B opdeles
på en samling af gensidigt udelukkende og udtømmende hændelser, er vist i \autoref{fig:decoAmulB}. Hændelsen A er opdelt i fem gensidigt udelukkende hændelser \( B_1, B_2, B_3, B_4, B_5 \), som tilsammen udgør udfaldsrummet \( S \). Den totale sandsynlighed for hændelse A beregnes som summen af de betingede sandsynligheder for A givet hver del \( B_i \) i partitionen ganget med sandsynligheden \( P(B_i) \) for hver del.

\begin{figure}[htbp]
    \centering
    \begin{tikzpicture}
        % Draw the rectangle to represent the universal set
        \draw[outline] (-2,-2) rectangle (8,4);
        
        % Draw the ellipse and fill
        \begin{scope}
            \fill[filled] \firstellipse;
        \end{scope}
        \draw[outline] 
            \firstellipse 
            node[anchor=north west, xshift=3.2cm, yshift=-1.5cm] {$A$};
        
        % Draw squiggly lines to create five partitions
        \draw[decorate, decoration={snake, amplitude=0.05cm, segment length=1cm, post length=0.5cm}] (0,4) -- (0,-2); % First squiggly line
        \draw[decorate, decoration={snake, amplitude=0.05cm, segment length=1cm, post length=0.5cm}] (2,4) -- (2,-2); % Second squiggly line (existing one)
        \draw[decorate, decoration={snake, amplitude=0.05cm, segment length=1cm, post length=0.5cm}] (4,4) -- (4,-2); % Third squiggly line
        \draw[decorate, decoration={snake, amplitude=0.05cm, segment length=1cm, post length=0.5cm}] (6,4) -- (6,-2); % Fourth squiggly line
        
        % Label the partitions B_1, B_2, ..., B_5
        \node at (-1, 3.5) {$B_1$};
        \node at (1, 3.5) {$B_2$};
        \node at (3, 3.5) {$B_3$};
        \node at (5, 3.5) {$B_4$};
        \node at (7, 3.5) {$B_5$};
    \end{tikzpicture}
    \caption{$P(A) = \sum_{i=1}^{5} P(A \mid B_i) P(B_i)$}
    \label{fig:decoAmulB}
\end{figure}

\begin{example}
    Et universitet udbyder tre typer kurser: \textbf{online}, \textbf{hybride} og \textbf{fremmødekurser}. Sandsynligheden for, at en studerende dumper et kursus, er \( P(F \mid \text{Online}) = 0.10 \) for onlinekurser, \( P(F \mid \text{Hybrid}) = 0.05 \) for hybride kurser og \( P(F \mid \text{Fremmøde}) = 0.02 \) for fremmødekurser. Universitetet udbyder 50\% af sine kurser online, 30\% som hybride kurser og 20\% som fremmødekurser. Hvad er sandsynligheden for, at en tilfældigt udvalgt studerende dumper et kursus?
\end{example}

\begin{solution}
    Med loven om total sandsynlighed får vi:

    \[
    P(F) = P(F \mid \text{Online}) P(\text{Online}) + P(F \mid \text{Hybrid}) P(\text{Hybrid}) + P(F \mid \text{Fremmøde}) P(\text{Fremmøde})
    \]

    Vi indsætter de givne værdier:

    \[
    P(F) = 0.10 \times 0.50 + 0.05 \times 0.30 + 0.02 \times 0.20 = 0.05 + 0.015 + 0.004 = 0.069
    \]

    Sandsynligheden for, at en tilfældigt udvalgt studerende dumper et kursus, er altså 0.069.
\end{solution}

\section{Uafhængighed}
Lad $A$ være hændelsen, at det regner i morgen, med $P(A)=\frac{1}{3}$. Antag desuden, at jeg kaster en fair mønt, og lad $B$ være hændelsen, at den lander med krone opad, så $P(B)=\frac{1}{2}$.

Se nu på sandsynligheden $P(A \mid B)$. Hvad tror du, den er? Du gætter måske intuitivt på, at $P(A \mid B)=P(A)=\frac{1}{3}$, og det er korrekt! Udfaldet af møntkastet har ingen indflydelse på vejret. Det betyder, at sandsynligheden for $A$ er uændret, uanset om $B$ indtræffer eller ej. Scenariet illustrerer begrebet uafhængige hændelser: To hændelser er uafhængige, hvis forekomsten af den ene ikke giver nogen information om den anden.

Lad os nu formalisere definitionen af uafhængighed.

\begin{definition}[Uafhængighed]
    To hændelser \( A \) og \( B \) er uafhængige, hvis
    \[
        P(A \cap B)=P(A)P(B).
    \]
    Når \(P(B)>0\), er dette ensbetydende med \(P(A\mid B)=P(A)\), og når \(P(A)>0\), er det ensbetydende med \(P(B\mid A)=P(B)\).
\end{definition}

Sammenfattende kan uafhængighed forstås på to ækvivalente måder: »Uafhængighed betyder, at sandsynligheden for fællesmængden af to hændelser kan findes ved blot at gange deres individuelle sandsynligheder,« eller alternativt: »Uafhængighed betyder, at den betingede sandsynlighed for den ene hændelse givet den anden er den samme som den oprindelige, ubetingede sandsynlighed.«

\begin{lemma}
Hvis $A$ og $B$ er uafhængige, så gælder:
\begin{itemize}
    \item $A$ og $B^c$ er uafhængige,
    \item $A^c$ og $B$ er uafhængige,
    \item $A^c$ og $B^c$ er uafhængige.
\end{itemize}
\end{lemma}

Når vi arbejder med sandsynligheden for foreningsmængden af flere uafhængige hændelser, \( A_1, A_2, \ldots, A_n \), er det ofte lettere at finde sandsynligheden for deres fællesmængde end for deres foreningsmængde. I sådanne situationer kan De Morgans lov være meget nyttig:

\[
A_1 \cup A_2 \cup \cdots \cup A_n = \left(A_1^c \cap A_2^c \cap \cdots \cap A_n^c\right)^c
\]

Ved hjælp af denne sammenhæng kan vi udtrykke sandsynligheden for foreningsmængden som:

\begin{align*}
P\left(A_1 \cup A_2 \cup \cdots \cup A_n\right) &= 1 - P\left(A_1^c \cap A_2^c \cap \cdots \cap A_n^c\right) \\
&= 1 - P\left(A_1^c\right) P\left(A_2^c\right) \cdots P\left(A_n^c\right) \quad \text{(da hændelserne \( A_i \) er uafhængige)} \\
&= 1 - \left(1 - P\left(A_1\right)\right)\left(1 - P\left(A_2\right)\right) \cdots \left(1 - P\left(A_n\right)\right).
\end{align*}


\begin{theorem}[Uafhængighed og De Morgans lov]
        Hvis \( A_1, A_2, \cdots, A_n \) er uafhængige, så gælder
        \begin{align*}
        &P\left(A_1 \cup A_2 \cup \cdots \cup A_n\right)=1-\left(1-P\left(A_1\right)\right)\left(1-P\left(A_2\right)\right) \cdots\left(1-P\left(A_n\right)\right)
        \end{align*}
\end{theorem}

\textbf{Advarsel!} En almindelig misforståelse er at forveksle uafhængighed med disjunkthed. Det er dog grundlæggende forskellige begreber. To hændelser $A$ og $B$ er disjunkte, hvis $A \cap B=\emptyset$. Hvis begge hændelser har positiv sandsynlighed, medfører disjunkthed afhængighed: Hvis vi ved, at den ene er indtruffet, er den anden udelukket. Undtagelsen er, når mindst én af de disjunkte hændelser har sandsynlighed nul. I det tilfælde er produktkriteriet for uafhængighed opfyldt.

\begin{table}[ht]
    \begin{center}
    \renewcommand{\arraystretch}{1.5} % Adjusts row spacing
    \begin{tabular}{c|c|c}
    \textbf{Begreb} & \textbf{Beskrivelse} & \textbf{Centrale formler} \\
    \hline
    Disjunkte & Hændelserne $A$ og $B$ kan ikke indtræffe samtidig & 
    \begin{tabular}{@{}c@{}}
    $A \cap B = \emptyset$ \\
    $P(A \cup B) = P(A) + P(B)$
    \end{tabular} \\
    \hline
    Uafhængige & Forekomsten af $B$ giver ingen information om $A$ &
    \begin{tabular}{@{}c@{}}
    $P(A \mid B) = P(A)$ \\
    $P(B \mid A) = P(B)$ \\
    $P(A \cap B) = P(A) \cdot P(B)$
    \end{tabular} \\
    \end{tabular}
    \end{center}
    \caption{Sammenligning af disjunkte og uafhængige hændelser.}
    \label{tab:disjoint_vs_independent}
\end{table}

Vi kan udvide begrebet uafhængighed til flere hændelser. En samling hændelser \( A_1, A_2, \ldots, A_n \) betragtes som uafhængige, hvis forekomsten af en vilkårlig delmængde af hændelserne ikke giver nogen information om forekomsten af de øvrige hændelser. Matematisk udtrykkes det således:

\begin{definition}[Uafhængighed af flere hændelser]
    Hændelserne \( A_1, A_2, \ldots, A_k \) er indbyrdes uafhængige, hvis det for enhver delmængde \(J\subseteq\{1,\ldots,k\}\) med mindst to indeks gælder, at
    \[
    P\left(\bigcap_{j\in J}A_j\right)=\prod_{j\in J}P(A_j).
    \]
    Det er ikke tilstrækkeligt kun at kræve lighedstegnet for fællesmængden af alle \(k\) hændelser.
\end{definition}

\begin{example} Spillerens fejlslutning

\textbf{Spillerens fejlslutning} er en almindelig kognitiv bias, som opstår, når man tror, at udfaldet af en tilfældig hændelse påvirkes af tidligere udfald. Fejlslutningen ses ofte i forbindelse med spil, hvor man fejlagtigt antager, at sandsynligheden for en hændelse påvirkes af tidligere hændelser.

En mand kaster en fair mønt otte gange og noterer ved hvert kast, om det giver krone $(\mathrm{H})$ eller plat $(\mathrm{T})$. Efter hvilken af følgende sekvenser af møntkast er det mest sandsynligt, at manden får krone $(H)$ i sit næste kast? Denne:

TTTTTTTT

eller denne:

HHTHTTHH

Svaret er ingen af dem, som det fremgår her:

\begin{align*}
        P\left(H_9 \mid T_1 T_2 \ldots T_8\right)=\frac{P\left(H_9 \cap T_1 \cap T_2 \cap \cdots \cap T_8\right)}{P\left(T_1 \cap T_2 \cap \cdots \cap T_8\right)} & =\frac{\left(\frac{1}{2}\right)^9}{\left(\frac{1}{2}\right)^8} \\
        & =\frac{1}{2}
\end{align*}
    
\end{example}

Prøv at undgå at falde i spillerens fejlslutning. Hvis en fair mønt for eksempel er landet på plat otte gange i træk, er det let at tro, at krone »står for tur« i næste kast. Men hvert kast er uafhængigt, og sandsynligheden er den samme -- der er stadig 50\% chance for krone eller plat, uanset tidligere resultater.

\section{Bayes' sætning}
Vi er nu klar til at introducere et af de stærkeste værktøjer inden for betinget sandsynlighed: Bayes' regel (eller sætning). Reglen er særlig nyttig, når vi kender \( P(A \mid B) \), men vil finde \( P(B \mid A) \). Ud fra definitionen af betinget sandsynlighed har vi:

\[
P(A \mid B) \cdot P(B) = P(A \cap B) = P(B \mid A) \cdot P(A)
\]

Ved at dividere begge sider med \( P(A) \) får vi:

\[
P(B \mid A) = \frac{P(A \mid B) \cdot P(B)}{P(A)}
\]

Denne formel er velkendt som Bayes' regel. I mange tilfælde har vi brug for loven om total sandsynlighed for at beregne \( P(A) \) i Bayes' regel. Derfor præsenteres Bayes' regel ofte på formen:

\[
P(B \mid A) = \frac{P(A \mid B) \cdot P(B)}{P(A \mid B) \cdot P(B) + P(A \mid B^c) \cdot P(B^c)}
\]

eller mere generelt

\[
P\left(B_j \mid A\right) = \frac{P\left(A \mid B_j\right) \cdot P\left(B_j\right)}{\sum_i P\left(A \mid B_i\right) \cdot P\left(B_i\right)}
\]

hvor \( B_1, B_2, \ldots, B_n \) udgør en partition af udfaldsrummet.

\begin{theorem}[Bayes' sætning]
    For to vilkårlige hændelser $A$ og $B$, hvor $P(A) \neq 0$, har vi
    \begin{align*}
    P(B \mid A) = \frac{P(A \mid B) \cdot P(B)}{P(A)}
     = \frac{P(A \mid B) \cdot P(B)}{P(A \mid B) \cdot P(B) + P(A \mid B^c) \cdot P(B^c)}
    \end{align*}

    Hvis $B_1, B_2, B_3, \cdots$ udgør en partition af udfaldsrummet $S$, og $A$ er en vilkårlig hændelse med $P(A) \neq 0$, har vi

    \begin{align*}
    P\left(B_j \mid A\right) = \frac{P\left(A \mid B_j\right) \cdot P\left(B_j\right)}{\sum_i P\left(A \mid B_i\right) \cdot P\left(B_i\right)}
    \end{align*}
\end{theorem}

Vi begynder med et berygtet eksempel, som fremhæver betydningen af Bayes' sætning.

\begin{example} Det falsk positive paradoks

   Forestil dig en sjælden sygdom, som rammer omkring 1 ud af hver 10\,000 personer. Der findes en test til at påvise sygdommen, og selv om testen er meget nøjagtig, er den ikke perfekt. Konkret gælder:
\begin{itemize}
    \item Sandsynligheden for, at testen viser et positivt resultat (som tyder på sygdom), når personen ikke har sygdommen, er $2 \%$.
    \item Sandsynligheden for, at testen viser et negativt resultat (som tyder på, at der ikke er sygdom), når personen har sygdommen, er $1 \%$.
\end{itemize}

Antag nu, at en tilfældigt udvalgt person bliver testet, og at resultatet er positivt. Hvad er sandsynligheden for, at personen faktisk har sygdommen?
\end{example}

\begin{solution}
    Lad \( D \) være hændelsen, at personen har sygdommen, lad \( T^+ \) være hændelsen, at testresultatet er positivt, og lad \( T^- \) være hændelsen, at det er negativt. Vi har fået oplyst:

\[
\begin{aligned}
& P(D) = \frac{1}{10\,000} \\
& P\left(T^+ \mid D^c\right) = 0.02 \quad \text{(Falsk positiv-rate)} \\
& P\left(T^- \mid D\right) = 0.01 \quad \text{(Falsk negativ-rate)}
\end{aligned}
\]

Vi vil beregne \( P(D \mid T^+) \). Først har vi brug for \( P(T^+ \mid D) \), som er sandsynligheden for et sandt positivt resultat. Da en person med sygdommen enten testes positiv eller negativ, har vi:

\[
P(T^+ \mid D) = 1 - P(T^- \mid D) = 1 - 0.01 = 0.99
\]

Med Bayes' regel får vi nu:

\[
\begin{aligned}
P(D \mid T^+) &= \frac{P(T^+ \mid D) P(D)}{P(T^+ \mid D) P(D) + P\left(T^+ \mid D^c\right) P\left(D^c\right)} \\
&= \frac{0.99 \times 0.0001}{0.99 \times 0.0001 + 0.02 \times (1 - 0.0001)} \\
&= \frac{0.000099}{0.000099 + 0.02 \times 0.9999} \\
&= \frac{0.000099}{0.000099 + 0.019998} \\
&\approx 0.0049
\end{aligned}
\]

Resultatet betyder, at der er mindre end en halv procents chance for, at personen faktisk har sygdommen. Trods det positive testresultat er det sygdommens lave forekomst og testens falsk positiv-rate, der fører til dette kontraintuitive resultat.
\end{solution}

\begin{example}
    Bayesianske netværk bruges ofte på højteknologiske producenters websteder for at hjælpe kunder med hurtigt at diagnosticere problemer med deres produkter. En printerproducent bruger for eksempel data fra testresultater til at finde mulige årsager til printerfejl. Printerfejl hænger primært sammen med tre typer problemer: hardware, software og andre problemer (som stik og forbindelser). Sandsynlighederne for disse problemer er som følger:
\begin{itemize}
    \item Sandsynligheden for et hardwareproblem: $P(H)=0.1$
    \item Sandsynligheden for et softwareproblem: $P(S)=0.6$
    \item Sandsynligheden for en anden type problem: $P(O)=0.3$
\end{itemize}

Sandsynligheden for, at en printer fejler, givet hver type problem, er:
\begin{itemize}
    \item Sandsynligheden for fejl givet et hardwareproblem: $P(F \mid H)=0.9$
    \item Sandsynligheden for fejl givet et softwareproblem: $P(F \mid S)=0.2$
    \item Sandsynligheden for fejl givet et andet problem: $P(F \mid O)=0.5$
\end{itemize}

Givet at en kunde oplever en printerfejl og bruger producentens websted til at diagnosticere problemet, hvad er så den mest sandsynlige årsag til problemet?
\end{example}

\begin{solution}
    For at bestemme den mest sandsynlige årsag til printerfejlen skal vi beregne sandsynligheden for hver type problem, givet at der er opstået en fejl. Det gør vi ved at bruge Bayes' sætning til at beregne a posteriori-sandsynlighederne.

Lad $F$ betegne hændelsen, at printeren fejler. Vi vil finde:

\[
P(H \mid F), \quad P(S \mid F), \quad P(O \mid F)
\]

\textbf{Trin 1: Beregn den totale sandsynlighed for fejl
}

Først bruger vi loven om total sandsynlighed til at finde $P(F)$:

\[
P(F) = P(F \mid H) \cdot P(H) + P(F \mid S) \cdot P(S) + P(F \mid O) \cdot P(O)
\]

Vi indsætter de givne værdier:

\[
\begin{aligned}
P(F) &= (0.9 \times 0.1) + (0.2 \times 0.6) + (0.5 \times 0.3) \\
P(F) &= 0.09 + 0.12 + 0.15 = 0.36
\end{aligned}
\]

\textbf{Trin 2: Brug Bayes' sætning til at finde a posteriori-sandsynlighederne
}

Brug nu Bayes' sætning for hver type problem:

\begin{enumerate}
    \item Sandsynligheden for et hardwareproblem givet fejl: \vspace{10pt}
    
    \[
    P(H \mid F)=\frac{P(F \mid H) \cdot P(H)}{P(F)}=\frac{0.9 \times 0.1}{0.36}=\frac{0.09}{0.36}=0.25
    \]
    \vspace{10pt}
    \item Sandsynligheden for et softwareproblem givet fejl: \vspace{10pt}
    
    \[
    P(S \mid F)=\frac{P(F \mid S) \cdot P(S)}{P(F)}=\frac{0.2 \times 0.6}{0.36}=\frac{0.12}{0.36}=0.3333
    \]
    \vspace{10pt}
    \item Sandsynligheden for andre problemer givet fejl: \vspace{10pt}
    
    \[
    P(O \mid F)=\frac{P(F \mid O) \cdot P(O)}{P(F)}=\frac{0.5 \times 0.3}{0.36}=\frac{0.15}{0.36}=0.4167
    \]
    \vspace{10pt}
\end{enumerate}

\textbf{Trin 3: Fortolk resultaterne}
\begin{itemize}
    \item Hardwareproblem: $25 \%$ chance
    \item Softwareproblem: cirka $33.33 \%$ chance
    \item Andet problem: cirka $41.67 \%$ chance
\end{itemize}

\textbf{Konklusion:} Givet en printerfejl er den mest sandsynlige årsag et andet problem, for eksempel stik og forbindelser, med en sandsynlighed på cirka $41.67 \%$.


\end{solution}
